% ORDER4GAP1
\documentclass[11pt]{article}
\usepackage[margin=.88in]{geometry}
\usepackage{amsmath,amssymb,amsthm}
\newtheorem{proposition}{Proposition}
\begin{document}
\begin{center}
{\Large ORDER4GAP1: after order three, the width-2 gap is an order-four sum}\\[4pt]
\textbf{Date: 2026-10-04 22:31 PT}\\
\textbf{Status: PROVED/EXACT bound + COMPUTATIONAL scale}\\
\textbf{P versus NP status: OPEN}
\end{center}

WIDTH2REM1 left the first unmatched inclusion at order four. This note bounds it.

\begin{proposition}[ORDER4-GAP]
Let $\mu$ and $\nu$ agree on every moment through order three on the same exact slice. For a width-2 monotone DNF with term set $\mathcal T$,
\[
\bigl|\mathbb E_\mu[\mathrm{DNF}]-\mathbb E_\nu[\mathrm{DNF}]\bigr|
\le
\sum_{\{T,T'\}\ \mathrm{disjoint}}
\Pr[T\cup T'\subseteq F],
\]
where the probability may be taken under either law only after replacing it by the exact-slice four-coordinate probability $(t)_4/(M)_4$. In particular the right-hand side is at most $\binom{|\mathcal T|}{2}(t)_4/(M)_4$.
\end{proposition}
\begin{proof}
The first Bonferroni sum is a sum of pair moments, hence equal under $\mu$ and $\nu$. Intersecting term pairs are three-coordinate moments, hence also equal. The truncation error of the union after these matched terms is at most the sum of the unmatched pairwise intersections, which are the disjoint term pairs. Each such intersection is a four-coordinate all-ones event, whose exact-slice probability is $(t)_4/(M)_4$.
\end{proof}

ALL-ONES-LOG at $r=4$ gives
\[
-\log\frac{(t)_4/(M)_4}{(t/M)^4}\le \frac{8}{t-3}.
\]
The numerical check holds: at $q=20$ the log-relative error is $0.0865$ against bound $0.119$.

The resulting scale is not small for a large DNF. At $q=20$ and $10{,}000$ terms it is $0.159$. The MOMENTBRIDGE1 mean gap $0.010$ sits inside it. A $1{,}000$-term DNF on the same slice has scale $0.00159$. So order-three matching kills the width-2 gap only for DNFs whose term count is $o(M^2/t^2)$, not for the probes already run.

The bound is generic to every fixed-weight slice, including the GE-easy XOR slice. It is not a monotone lower bound. $P\ne NP$ remains OPEN.
\end{document}
